\chapter{Conclusions}
\label{ch:conclusion}

This chapter draws together the empirical work of the thesis and sets it against the broader questions that motivated it. Section~\ref{sec:contribution-summary} summarises the contributions across the three research questions. Section~\ref{sec:limitations} states the limitations honestly, including the ones we discovered along the way. Section~\ref{sec:future-directions} outlines the extensions that follow most naturally from the current results. Section~\ref{sec:closing-thoughts} closes with a short reflection on what the thesis, taken as a whole, suggests about the study of emergent communication.

\section{Summary of Contributions}
\label{sec:contribution-summary}

The thesis makes four main contributions.

\paragraph{An empirical framework for studying emergent language under structured input.}
The thesis extended the PACE architecture to accept Zipfian training distributions parametrised by a concentration parameter $\alpha$, ran the resulting training campaign at scale, and built the analysis pipeline needed to compare populations across seeds and concentration levels. All experiments were executed on the Vera cluster at Chalmers, and the codebase, the training checkpoints, and the analysis scripts are available in the project repository for reproducibility. The infrastructure was designed to keep the training procedure fixed and to vary only the input distribution, so that any observed effect could be attributed to concentration or ranking rather than to hyperparameter tuning.

\paragraph{A three-regime characterisation of how concentration shapes the internal language.}
Chapter~\ref{ch:experimental} showed that concentration shapes the internal language in a non-monotone way. Moderate $\alpha$ produces the largest and most compressed libraries, while both extremes reduce the resulting library: near-uniform input weakly differentiates among candidates, and highly concentrated input removes too much candidate diversity. A complementary analysis showed that in 566 of the 724 adoption steps (78.2\%), the adopted candidate differed from the candidate with the highest raw frequency, so the emergent library is organised around compressive value rather than raw exposure alone.

\paragraph{A quantitative demonstration of dialect formation across populations.}
Chapter~\ref{ch:rq2} showed that populations sharing the same concentration profile but differing in ranking develop systematically different abstraction libraries, and that the divergence increases monotonically with concentration. At low $\alpha$, pairwise Jaccard similarity is around 0.30; at $\alpha = 0.00$, it reaches approximately 0.34 and exceeds the corresponding null expectation by about 0.22. A null model based on random sampling from the global pool of adopted abstractions distinguished two regimes. In the low and moderate concentration regime, populations share more abstractions than random sampling would predict, a positive dialect signal indicating convergence on a shared structural backbone. In the high concentration regime, the observed overlap falls below the null, indicating that libraries become tightly tied to seed-specific top pairs and diverge from each other more than random sampling would produce. A qualitative analysis at moderate concentration showed that the divergence is not purely stochastic: the early abstractions adopted by each seed appear as sub-structures of that seed's top-ranked pairs, which mechanistically links the training distribution to the resulting library.

\paragraph{A structural picture of cross-population robustness.}
The picture that emerges from RQ3 is that cross-population communication in PACE is not primarily limited by neural adaptation but by structural properties of the training conditions. Chapter~\ref{ch:rq3} tested how the population-level differences shown in Chapter~\ref{ch:rq2} translate into communication difficulties when agents from different populations are paired. The chapter organised its experiments into same-$\alpha$ cross-pairing (C), mechanistic modifications of that setup (M), cross-$\alpha$ pairing under two speaker modes (X), and case studies further from the training conditions (O). In the same-$\alpha$ setting, adaptation difficulty depends strongly on the shared concentration regime and only weakly on pairwise distributional distance. In the cross-$\alpha$ setting, concentrated speakers produce libraries too specialised to pair robustly with listeners from other regimes, while uniform speakers generalise more broadly. Allowing the speaker to adapt neurally during joint training does not close this gap; the limiting factor is the discrete vocabulary built during independent training, not a neural mismatch that gradient descent can smooth out.

\section{Limitations}
\label{sec:limitations}

The results of this thesis are bounded by several limitations, some structural to the setup and some contingent on the choices we made.

\paragraph{Scale of the cross-population and modification analyses.}
The Chapter~\ref{ch:rq3} modifications (M1, M2, and M3) were tested on six listener seeds for M1 and M2, and on two listener seeds for M3. This pilot scale is sufficient to show the qualitative patterns reported in Sections~\ref{sec:rq3-m1} through~\ref{sec:rq3-m3}, but it does not support strong statistical claims. The cross-$\alpha$ analyses (X1 and X2) were likewise run for two speaker seeds, 42 and 123. The results replicate qualitatively across the two seeds, supporting the robustness of the observed asymmetry, but a fuller campaign would rotate the speaker across all twenty seeds and report the cross-pairing outcomes as an average over speaker identity.

\paragraph{Single reference seed for the case studies.}
The O2 and O3 case studies were run at a single seed, seed 42. The three sub-dynamics identified qualitatively at the $1.80 \rightarrow 2.40$ pairing (substitution, abandonment, and emergence) may reflect properties specific to that seed's library rather than general properties of PACE. Replicating on other seeds would clarify whether the substitution dynamic in particular is a canonical phenomenon or an idiosyncrasy of this specific library.

\paragraph{Fixed training horizon.}
All PACE runs in this thesis inherit the training horizon of the original architecture, which is 19 abstraction proposal steps. In the intermediate concentration regime, the candidate pool remains rich enough that new abstractions might continue to emerge if training were extended, and the library sizes reported in Chapter~\ref{ch:experimental} could in principle be lower bounds. Extending the horizon would clarify whether the reported results capture the full dynamics of lexical innovation or whether the system is still discovering when the campaign ends.

\paragraph{Usage of the emergent vocabulary.}
Our instrumentation measures how many abstractions each population adopts and how much they compress the message stream, but not how evenly those abstractions are used across the training scenes. Whether a concentrated input produces a vocabulary that is merely smaller, or one that is also used more unequally, is a distinct question that the current logging does not answer and that would require per-abstraction usage counts at convergence.

\paragraph{Scope of the environment.}
The interpretation of the results is bounded by the specific environment that PACE assumes: a fixed scene inventory built from a small set of primitive drawing actions, a grid-based world, and the specific combination of deterministic MDL-based abstraction adoption with $\varepsilon$-greedy program selection. It also deliberately excludes large-scale visual or linguistic pretraining. This makes the origin of PACE's abstractions easier to isolate, but separates the setup from multimodal systems that transfer knowledge across language and perception~\cite{huang_language_2023}, and from vision-language-action models that use web-scale pretraining to generalise robotic control to novel objects and instructions~\cite{brohan_rt-2_2023}. Whether the concentration effects found here persist in such richer, prior-informed settings remains an open question.

\section{Future Directions}
\label{sec:future-directions}

Several extensions follow naturally from the current results.

\paragraph{Scaling up the modifications and cross-population analyses.}
The most immediate extension is to increase the sample size in the analyses currently run at pilot scale. Running M1, M2, and M3 across the full 20-seed pool would let us test whether the qualitative patterns identified in Chapter~\ref{ch:rq3} hold at population level. Running the full cross-$\alpha$ grid across all 20 speaker seeds, rather than the two currently used, would let us report the cross-pairing outcomes as an average over speaker identity and remove the dependence on the specific libraries of seeds 42 and 123.

\paragraph{Rotating the reference agent.}
Related to the previous point, the same-$\alpha$ experiments of Section~\ref{sec:rq3-group-c} use seed 42 as the fixed speaker in all cross-pairings. Rotating the speaker role across seeds would let us report a symmetric cross-pairing matrix and check whether the seed-42 speaker's library happens to be more or less friendly to cross-population adaptation than a typical speaker's library.

\paragraph{Testing extended concentration ranges.}
The current campaign spans $\alpha$ from 0.00 to 3.00, a range chosen because at $\alpha = 3.00$ the effective vocabulary already collapses to a handful of pairs. Extending the range downward, into fractional concentration values, or upward into a regime where all training mass concentrates on a single pair, would clarify what happens at the boundaries of the setup. Extension to negative values, in which the distribution is reversed and initially rare items become the most frequent, would test whether the observed patterns are specific to Zipfian shape or generalise to other skewed profiles.

\paragraph{Alternative bandit configurations and interaction structures.}
PACE in its current form uses $\varepsilon$-greedy program selection. The modification M3 in Chapter~\ref{ch:rq3} already shows that the exploration parameter has a measurable effect on the bandit's Q-value dynamics under cross-pairing. A fuller analysis would test how the adaptation results generalise to different bandit strategies or, more broadly, to setups in which PACE trains multiple architect-builder pairs simultaneously in a network structure, connecting PACE to population-level emergent-communication work~\cite{kim2021emergent,rita_role_2022}. Such a setup would let us probe how emerging conventions propagate across a population of interacting agents rather than emerging within a single dyad.

\paragraph{Generalisation to unseen scenes.}
A further axis, orthogonal to the ones studied here, is generalisation to unseen scenes. In human dyads, ad hoc conventions extend to undiscussed referents in a graded way, decaying with visual distance from the discussed targets~\cite{ji_ad_2025}. A complementary human result comes from geometric puzzle solving: He et al.~\cite{he_bootstrapping_2025} show that a solution schema abstracted from simple shape-construction puzzles can transfer to compositionally harder ones, including far transfer without surface similarity. Whether PACE abstractions support the analogous compositional generalisation to structurally related but unseen scenes, and whether that generalisation decays with a distance metric on the scene space, is a natural extension of RQ1 that the current fixed-inventory setup does not address.

\paragraph{Networked population structure.}
A further extension concerns population structure. Here, populations are trained in isolation and paired only at evaluation time. In human cultural evolution, the topology of the social network through which artefacts propagate measurably shapes what emerges~\cite{marjieh_characterizing_2025}. Embedding multiple PACE populations in a network and letting abstractions propagate between neighbours would let one ask how topology mediates the convergence or fragmentation of emergent vocabularies, an axis orthogonal to the input-concentration effects studied here. Beyond a fixed topology, the interaction structure could itself be adaptive: Leung et al.~\cite{leung_curiosity-driven_2025} show in an emergent-communication setting that partner selection reshapes who interacts with whom and can accelerate convergence to shared conventions. A PACE version of this setup would ask whether abstraction libraries converge or fragment depending on which populations end up communicating most often.

\paragraph{Cross-population communication beyond joint training.}
All the cross-pairing experiments in this thesis measure adaptation, that is, how quickly the two agents come to understand each other during a joint training phase. A complementary measurement would be intelligibility at epoch zero, before any joint training, which isolates the raw compatibility of the two agents' vocabularies from the effect of gradient descent. The X1 and X2 sweeps of Chapter~\ref{ch:rq3} already report epoch-0 accuracy as one of their aggregate metrics, but a systematic analysis of what predicts epoch-0 intelligibility across the cross-$\alpha$ grid would provide a cleaner picture of the structural gap identified in Section~\ref{sec:rq3-group-x}.

\section{Closing Thoughts}
\label{sec:closing-thoughts}

The work developed in this thesis points to three broader lessons for emergent communication. Input shape matters, but concentration does not act like a single dial that smoothly shrinks or specialises the language. Its effect depends on the balance between repetition and diversity, and that balance changes qualitatively across the range. Populations trained under the same concentration are also neither fully aligned nor fully independent: they share some structural pressure, but the specific ranking of frequent scenes still leaves a strong imprint.

This convergence-and-diversity trade-off echoes a finding from a purely neural setting: Sevestre and Dupoux~\cite{sevestre_frequency_2025} argue that it is not frequency per se but the limited data exposure it creates that drives structural change in emergent languages, a framing continuous with iterated-learning accounts of how bottlenecks promote structure~\cite{kirby_spontaneous_2001}. Their result concerns compositionality over frequent and rare objects, whereas ours concerns the size and transferability of an adopted symbolic library. The two findings are therefore complementary rather than redundant. Our three-regime picture is the symbolic-vocabulary counterpart: concentration helps only up to the point where it starves the diversity the abstraction mechanism needs.

A second lesson concerns the relationship between PACE's symbolic and neural components. The symbolic side (which abstractions are proposed, adopted, and retained) and the neural side (how quickly the agents learn to use them) are tightly coupled without being interchangeable. Chapter~\ref{ch:rq3} gives the clearest example: the cross-$\alpha$ gap survives even when both agents co-adapt neurally during joint training. Fast neural adjustment can reduce local difficulty, but it cannot repair a mismatch in the discrete vocabularies built during independent training.

One methodological lesson from this work is that inspectable neuro-symbolic systems like PACE allow questions about the internal structure of emergent languages to be posed at a level of granularity that is difficult to reach with purely neural approaches. Because PACE keeps its vocabulary inspectable, the analysis can move from aggregate accuracy curves to concrete symbolic outcomes: which abstractions were adopted, which remained active during adaptation, and which were abandoned or replaced. The value of the symbolic component here is therefore not only computational; it also makes the scientific object easier to inspect.

The study of emergent communication sits at an unusual intersection. It is a subfield of machine learning, but its object of study, how communication systems form and stabilise, is one of the oldest and richest questions in linguistics and cognitive science. The concepts and methods we bring to bear on it come from both traditions, and the phenomena we discover often reflect the mechanistic details of the specific system we are studying more than the general principles we hoped to isolate. The work presented here has aimed to be specific about the mechanisms and honest about the reach of its findings. If it leaves the reader with the sense that the concentration of an input distribution shapes the emergent language in patterned but non-trivial ways, and that cross-population communication rests on more than pairwise similarity between training conditions, it will have done what it set out to do.